\newpage
\section{Acta Número 07}
\label{sec:acta07}

\noindent \textbf{Datos de la Sesión}
\begin{itemize}[noitemsep]
    \item \textbf{Modalidad:} Virtual - Google Meet.
    \item \textbf{Hora de Inicio:} 18:00
    \item \textbf{Hora de Fin:} 20:00
\end{itemize}

\noindent \textbf{Orden del Día}
\begin{itemize}[noitemsep]
    \item \textbf{Asistencia.}
    \item \textbf{Coordinación de Mockups Preliminares:} Planificación del diseño de interfaces a nivel general para las historias asignadas, previo al diseño responsivo detallado del Sprint 1.
    \item \textbf{Esquematización de la Base de Datos:} Coordinación y diseño inicial del modelo entidad-relación y esquemas de datos.
    \item \textbf{Firma y aprobación de los presentes.}
\end{itemize}

\noindent \textbf{Desarrolladores (Socios Fundadores)}
\begin{enumerate}[noitemsep]
    \item Pardo Pozo Juan Diego (Jefe de Proyecto).
    \item Arnez Jaimes Mauricio.
    \item Quispe Flores Camila Belen.
    \item Ariñez Bautista Jim Gabriel.
    \item Medina Rodríguez Mateo Jhoustin.
    \item Molina Maldonado Tomas Manuel.
    \item Zeballos Crespo Jonathan.
    \item Gutierrez Guzmán Angheline.
\end{enumerate}

\noindent \textbf{Fecha de la sesión:}\par
Jueves 26 de marzo de 2026

\newpage
\subsection{Desarrollo del Acta}

\noindent \textbf{Asistencia}\par
Se procedió a la toma de asistencia a la hora acordada, corroborando la presencia en sesión de los 8 socios fundadores de la grupo-empresa Byte Busters S.R.L. Al contar con la totalidad de los miembros, se estableció el quórum necesario para iniciar la toma de decisiones.

\vspace*{0.5cm}
\noindent\textbf{Coordinación de Mockups Preliminares}\par
El equipo coordinó la creación de wireframes y mockups de baja fidelidad para las Historias de Usuario asignadas recientemente. Se estableció que esta fase tiene como objetivo principal validar la estructura visual, la distribución de la información y el flujo de navegación básico. Se acordó explícitamente que el diseño detallado de la interfaz de usuario (UI), la aplicación de la guía de estilos y su respectiva adaptabilidad a dispositivos móviles (diseño responsivo) serán ejecutados y refinados a profundidad durante el desarrollo del Sprint 1.

\vspace*{0.5cm}
\noindent \textbf{Esquematización de la Base de Datos}\par
Se debatió y planificó la estructura de persistencia del sistema. El equipo comenzó a trazar los esquemas iniciales de la base de datos en PostgreSQL, identificando las entidades principales, sus atributos clave y las relaciones necesarias para soportar las funcionalidades del proyecto. Se asignaron responsabilidades para digitalizar el Modelo Entidad-Relación (MER) y definir la normalización de las tablas antes de proceder con los scripts de creación.

\vspace*{0.5cm}
\noindent \textbf{Aprobación Oficial}\par
Habiendo estructurado el enfoque visual preliminar y las bases del modelo de datos, los 8 socios fundadores aprobaron por unanimidad las decisiones tomadas durante la sesión, consolidando la dirección técnica previa al inicio del Sprint 1.

\newpage
\textbf{Firmas y aprobación de los integrantes}
\begin{center}
    \noindent
    \setlength{\tabcolsep}{0pt}
    
    \begin{tabular}{ccc}
        \espacioFirma{assets/img/firmas/mauricio.jpg}{Mauricio Jaimes Arnez}{13751221} & 
        \espacioFirma{assets/img/firmas/camila.jpg}{Camila Belen Quispe Flores}{13097244} &
        \espacioFirma{assets/img/firmas/jim.jpg}{Jim Gabriel Ariñez Bautista}{9517642} \\[1.0cm]
        
        \espacioFirma{assets/img/firmas/mateo.jpg}{Mateo Medina Rodríguez}{9453809} &
        \espacioFirma{assets/img/firmas/tomas.jpg}{Tomas Molina Maldonado}{8051701} &
        \espacioFirma{assets/img/firmas/jonathan.jpg}{Jonathan Zeballos Crespo}{13067494} \\[1.0cm]
    \end{tabular}

    \vspace{0.2cm} 
    \begin{tabular}{cc}
        \espacioFirma{assets/img/firmas/diego.jpg}{Juan Diego Pardo Pozo}{12747783} & 
        \espacioFirma{assets/img/firmas/angheline.jpg}{Angheline Gutierrez Guzmán}{13751815}
    \end{tabular}
\end{center}

\vspace{0.5cm}
\noindent \textbf{Fecha de última revisión:} 30 de marzo de 2026